\section{Product Features}

The following features describe externally visible capabilities of the initial
product. They express user value and safety behavior rather than prescribing a
specific implementation.

\subsection{PF-01: Multimodal Medication and Request Intake}

Users can upload or capture a printed/digital Vietnamese prescription, enter
medicine information manually, type a request, or speak a request. OCR converts
prescription content to editable fields, while speech-to-text converts voice to
an editable transcript. The interface shows recording/processing state,
permission errors, extraction confidence or uncertainty, and a manual fallback.
This feature is high priority.

\subsection{PF-02: Verification and Correction}

Before medication identification or assessment, users can inspect and correct
extracted or transcribed medicine names, strengths, dosage instructions, and
schedules. Low-confidence, missing, conflicting, and unresolved fields are
visibly identified. Explicit confirmation creates the verified input used by
downstream functions. The user confirms that the extracted details match the
source, not that they are medically correct. Any later edit invalidates the
previous normalization, mapping, and assessment and causes them to be run again.
This feature is high priority and safety-critical.

\subsection{PF-03: Intent Identification and Conversational Assistance}

The system identifies one or more supported intents from typed or confirmed
transcribed requests, such as checking interactions, explaining a structured
finding, managing a reminder, or recording a possible reaction. It requests
clarification when the intent or medicine is ambiguous. Conversational output
may restate structured findings in plain language but cannot create an
independent safety conclusion. This feature is high priority.

\subsection{PF-04: Vietnamese Medicine Identification and Normalization}

The system matches verified Vietnamese medicine information to prepared DAV
records, identifies active ingredients and relevant registration attributes,
and preserves source and version information. Ambiguous or unmatched products
are clearly unresolved. Prepared data and active-ingredient mappings are
versioned and reviewed before use. This feature is high priority and
safety-critical.

\subsection{PF-05: Medication Safety and Interaction Assessment}

Reliably identified active ingredients are mapped to reviewed DDInter 2.0
entities. The system checks supported drug--drug interactions and severity. A
result has an explicit state: interaction detected, no known interaction found
within the available data, incomplete, unassessed, or temporarily unavailable.
Before assessment, the system checks for required context such as the complete
medicine list, dose, route, allergies, age, pregnancy, symptoms, and relevant
kidney or liver impairment. It requests missing information instead of assuming
that an unknown condition is safe. The resulting route is routine/no significant
risk detected from available information, professional review required or
uncertain, or emergency/red flags.
Therapeutic duplication and drug--disease checks are included only when the
team can establish adequate data, mapping, and acceptance tests. This feature
is high priority and safety-critical.

\subsection{PF-06: Safety Findings and Explainable Guidance}

Users receive the involved medicines, severity where available, a concise
explanation grounded in the structured interaction record, source/coverage
limitations, and an appropriate next step. Warnings never instruct users to
start, stop, replace, or change a prescribed medicine independently. Severe,
uncertain, unresolved, or concerning cases advise qualified professional
review. Emergency red flags receive immediate instructions to seek urgent
real-world help and never wait for the simulated review workflow. This feature
is high priority and safety-critical.

\subsection{PF-07: Professional Review Support}

The prototype prepares a structured case summary containing verified medicines,
active ingredients, mapping status, detected findings, limitations, schedules,
and reported symptoms. A simulated professional reviewer can record guidance or
schedule confirmation, or request more information from the patient before the
case is reassessed. A case without a response remains pending or unavailable;
the system provides conservative interim guidance and recommends an external
doctor or pharmacist rather than implying approval. The product labels this
workflow as a simulation and does not imply that real clinical review has
occurred. This feature is high priority for demonstrating the safety workflow.

\subsection{PF-08: Confirmed-Schedule Medication Reminders}

After reviewing and confirming the medication plan, users can separately opt in
to reminders based on verified prescription instructions and choose notification
times that remain within those instructions. Plan confirmation does not
automatically activate reminders. Changes to the prescribed dosing schedule
require professional confirmation and user acceptance. A new reminder for a
medicine involved in a detected severe interaction remains on hold until review
and confirmation; unrelated reminders may proceed. This feature is medium
priority.

\subsection{PF-09: Adherence and Possible Adverse-Reaction Follow-up}

Users can respond to a reminder by recording a dose as taken, missed, or skipped,
and the system stores the resulting adherence event. Users can also record
symptoms for follow-up. Symptoms are labeled as possible adverse drug reactions,
not proven causation, and are triaged into mild/self-monitoring, professional
review required, or emergency/red flags. Review and emergency outcomes follow
their respective routes. This feature is medium priority.

\subsection{PF-10: Medication and Assessment History}

Users can review their verified medicine list, prior assessments, result states,
reminders, adherence events, possible-reaction records, and review guidance.
History preserves the relevant data-source/version context so that a previous
result is not presented as a current guarantee. This feature is medium priority
and may be reduced if the prototype cannot provide adequate privacy controls.

